\documentclass[11pt]{article}
\usepackage[a4paper,margin=2.5cm]{geometry}
\begin{document}
\begin{center}
{\Large \bf Report on ``Supersymmetric solutions of non-relativistic 11-dimensional supergravity''}\\[4pt]
{\large (C.~D.~A.~Blair, IFT-UAM/CSIC-25-26)}
\end{center}

\section*{Summary of the paper}
The paper studies the membrane (M2) non-relativistic limit of eleven-dimensional supergravity (``M2NC SUGRA''). This theory comes in two versions: a maximally supersymmetric one with a heavy tower of constraints, and a ``half-maximal'' one with fewer constraints and $\epsilon_-=0$. The author takes limits of the M2 and M5 brane solutions in two ways. In the \emph{scaled} limit the brane number is rescaled by hand, and in the \emph{smeared} limit the branes are smeared over the longitudinal directions. The paper checks the equations of motion, the constraints, and the Killing spinor equations for each limit. It also notes that longitudinal AdS$_3$ and the BTZ black hole solve the theory.

\section{Principal achievement}
The paper's main achievement is a clean and explicit dictionary between the two ways of obtaining non-Lorentzian brane backgrounds and the two versions of the supergravity. Scaled, longitudinally localised solutions satisfy the maximally supersymmetric constraints. Smeared, transversely localised solutions satisfy only the half-maximal constraints. The solutions are shown to be supersymmetric by solving the Killing spinor equations explicitly, and several infinite-dimensional supersymmetry enhancements are identified. This gives a non-trivial consistency check on the constraint structure of Bergshoeff et al., which was partly based on linearised arguments. It also gives a reason to take the unconventional half-maximal version seriously for non-relativistic holography. The author also argues that the pattern should hold for the other string and brane limits.

\section{Principal drawback}
The main drawback is that the results are verified only on a narrow set of highly symmetric, effectively trivial backgrounds. The scaled solutions satisfy the equations of motion trivially, since the longitudinal dependence is unconstrained. The paper itself concedes that the physical meaning of the extra scaling is puzzling. The half-maximal theory has supersymmetry transformations that do not close into diffeomorphisms, and there is no criterion for counting preserved supercharges once infinite-dimensional enhancements appear. The whole analysis also assumes that the constraints and transformations of the earlier work are correct, and the author acknowledges that this earlier work rests partly on linearised evidence. Finally, the conclusions about other limits are conjectural, and the brane sources and the role of the scaling of $N$ are left open. The claims are therefore consistency statements rather than a derivation of the physical significance of either limit.

\section{Most promising follow-up idea}
Balancing feasibility and impact, I would pursue coupling the solutions to worldvolume sources in the simplest dimensionally reduced setting. This is the check the author proposes: the M2NC limit of the M2 brane, reduced on a worldvolume direction, gives a string in the D2 Newton--Cartan limit (Matrix 2-brane theory). The supergravity limit then follows by dimensional reduction, and the string worldsheet action is already known. The coupled equations $S_{\rm SUGRA}+N S_{\rm worldvolume}$ can then be checked directly. This is feasible because every ingredient already exists, and it is a concrete calculation rather than a programme. It would also be high-impact, because it tests whether the smeared solutions, with no extra scaling of $N$, are genuinely sourced by the limit of the brane action. This addresses the main drawback above. It would decide between the scaled and smeared proposals as holographic backgrounds, and it would fix which version of the supergravity (maximal or half-maximal) is physically relevant. A natural extension is a separate check of whether the dualised null-compactified (DLCQ) frames remain consistent with the sourced solution.

\end{document}
